% Three ways to build a modern legal AI system (Buffa / Ren / Hamza). Vector TikZ.
% Design goal: three parallel pipelines, short labels, one highlighted step each.
% Full page width via figure* + resizebox.
\begin{figure*}[t]
\centering
\resizebox{\textwidth}{!}{%
\begin{tikzpicture}[
  every node/.style={font=\footnotesize},
  hdr/.style={rounded corners=2pt, align=center, text width=4.2cm, inner sep=4pt,
              font=\footnotesize\bfseries, text=white, minimum height=9mm},
  stp/.style={draw, rounded corners=2pt, align=center, text width=4.2cm,
               minimum height=8mm, inner sep=3pt, fill=white},
  key/.style={draw, line width=1pt, rounded corners=2pt, align=center, text width=4.2cm,
              minimum height=8.5mm, inner sep=3pt, font=\footnotesize\bfseries},
  tag/.style={align=center, text width=4.4cm, font=\scriptsize\itshape, text=black!55},
  fl/.style={-{Latex[length=2.2mm]}, draw=black!45, line width=0.7pt}]

  \definecolor{cBuffa}{HTML}{3B6EA5} % blue
  \definecolor{cRen}{HTML}{D98E2B}   % amber
  \definecolor{cHamza}{HTML}{BE3B39} % crimson

  \def\cA{0}     % column x-centres (cm)
  \def\cB{5.6}
  \def\cC{11.2}
  \def\rowH{1.2}

  % ===== Column 1 : Buffa et al. =====
  \node[hdr, fill=cBuffa] (b0) at (\cA,0) {Buffa et al.\\\normalfont\scriptsize RAG document building};
  \node[stp, draw=cBuffa] (b1) at (\cA,-\rowH)   {User query};
  \node[stp, draw=cBuffa] (b2) at (\cA,-2*\rowH) {Retrieve sources};
  \node[key, draw=cBuffa, fill=cBuffa!12, text=cBuffa!55!black] (b3) at (\cA,-3*\rowH) {Generate draft};
  \node[stp, draw=cBuffa] (b4) at (\cA,-4*\rowH) {Legal document};
  \node[tag] at (\cA,-4.95*\rowH) {Adapt the \textbf{generation}};
  \foreach \a/\b in {b0/b1,b1/b2,b2/b3,b3/b4} \draw[fl] (\a)--(\b);

  % ===== Column 2 : Ren et al. =====
  \node[hdr, fill=cRen] (r0) at (\cB,0) {Ren et al.\\\normalfont\scriptsize Fine-tuned legal RAG};
  \node[stp, draw=cRen] (r1) at (\cB,-\rowH)   {Synthetic legal data};
  \node[key, draw=cRen, fill=cRen!15, text=cRen!55!black] (r2) at (\cB,-2*\rowH) {Fine-tune retriever + generator};
  \node[stp, draw=cRen] (r3) at (\cB,-3*\rowH) {Retrieve + generate};
  \node[stp, draw=cRen] (r4) at (\cB,-4*\rowH) {Answer};
  \node[tag] at (\cB,-4.95*\rowH) {Adapt the \textbf{model}};
  \foreach \a/\b in {r0/r1,r1/r2,r2/r3,r3/r4} \draw[fl] (\a)--(\b);

  % ===== Column 3 : Hamza et al. =====
  \node[hdr, fill=cHamza] (h0) at (\cC,0) {Hamza et al.\\\normalfont\scriptsize Legalyze (SAHI)};
  \node[stp, draw=cHamza] (h1) at (\cC,-\rowH)   {Statute hierarchy};
  \node[key, draw=cHamza, fill=cHamza!12, text=cHamza!60!black] (h2) at (\cC,-2*\rowH) {SAHI index + breadcrumb};
  \node[stp, draw=cHamza] (h3) at (\cC,-3*\rowH) {Dual-stream + RRF};
  \node[stp, draw=cHamza] (h4) at (\cC,-4*\rowH) {Cited, addressed answer};
  \node[tag] at (\cC,-4.95*\rowH) {Adapt the \textbf{index}};
  \foreach \a/\b in {h0/h1,h1/h2,h2/h3,h3/h4} \draw[fl] (\a)--(\b);

\end{tikzpicture}}
\caption{Three ways to build a modern legal system, drawn as parallel pipelines. \citet{buffa2025}
adapt the \emph{generation} side: retrieve supporting text and let a model assemble a draft.
\citet{ren2026} go further and adapt the \emph{model} itself, fine-tuning both the retriever and the
generator on synthetic legal data. \citet{hamza2026} instead adapt the \emph{index}: they preserve
the statute's hierarchy with a breadcrumb per chunk, retrieve through two streams fused by rank, and
return each passage with its exact address. The highlighted box in each column marks that system's
defining step. The three answer different questions, and a mature system is likely to combine them.}
\label{fig:philosophies}
\end{figure*}
